\section{INTRODUCCIÓN}
En esta práctica se desarrolló la simulación electromagnética de un combinador de guía de onda WR42 operando alrededor de los 20 GHz, con el propósito de evaluar su comportamiento como estructura de combinación de potencia en sistemas de microondas. A diferencia de un análisis puramente teórico, el trabajo se centró en reproducir el proceso completo de diseño en HFSS: construcción geométrica del dispositivo, definición de condiciones de frontera, asignación de puertos de excitación, configuración de la solución y posterior interpretación de los resultados obtenidos.El dispositivo estudiado corresponde a una unión de cuatro puertos diseñada para combinar señales provenientes de dos entradas manteniendo bajas pérdidas y un adecuado aislamiento entre trayectorias. En este tipo de estructura, el comportamiento no depende únicamente de la forma física de la guía, sino también de la relación de fase entre las señales aplicadas, ya que una diferencia de fase apropiada permite que la potencia se sume de manera constructiva hacia el puerto de salida y se reduzca la energía dirigida al puerto aislado.

Durante el laboratorio, la simulación permitió observar cómo las decisiones de modelado influyen directamente en la respuesta electromagnética del sistema. El uso de simetría, la correcta asignación de fronteras conductoras y la definición adecuada de los puertos fueron aspectos fundamentales para obtener resultados confiables sin aumentar innecesariamente el costo computacional. Posteriormente, mediante los parámetros de dispersión y la visualización del campo eléctrico, se evaluó la adaptación del dispositivo, la transferencia de potencia entre puertos y la distribución de energía dentro de la guía de onda.Este informe presenta el procedimiento seguido y el análisis de los resultados principales, con énfasis en la interpretación física de la simulación. De esta manera, la práctica no se limita a comprobar el funcionamiento de una herramienta computacional, sino que permite relacionar el modelo numérico con el comportamiento real esperado de un combinador de potencia en alta frecuencia.
